% संपूर्ण मराठी ग्रंथातील प्रकरण 70: नैसर्गिक निगमन
% हा संपादनयोग्य स्रोतखंड आहे. संकलनासाठी संपूर्ण स्रोतसंग्रहातील
% अक्षररूपे, पूर्वभाग, चिन्हव्याख्या आणि आकृती-स्रोत आवश्यक आहेत.
% मूळ: Open Logic Project; मजकूर आणि रूपांतर CC BY 4.0.
% यंत्रानुवाद, दुरुस्ती आणि तपासणी: OpenAI Codex —
% GPT-5.6 Sol आणि GPT-6 Sol, दोन्ही Ultra effort.
% दुरुस्ती-पुनरावलोकन: GPT-6.1 Sol, Ultra effort.
\chapter{नैसर्गिक निगमन}\label{pt:nat:chap}









\section{परिचय}\label{pt:nat:int:sec}

नैसर्गिक निगमन हा सिद्धता
पद्धतींचा एक परिवार आहे.
त्यात प्रत्येक तार्किक संयोजकासाठी
किंवा संख्यापकासाठी दोन
अनुमाने असतात: एक संकारकाचे
“प्रवेशन” करणारे
(परिचय नियम)
आणि दुसरे त्याचे “विलोपन”
करणारे (निरसन नियम).
उदाहरणार्थ, \Intro{\lif}
च्या निष्कर्षात
$A \lif B$ या रूपाची
सूत्रे असतात,
तर \Elim{\lif} च्या
आधारविधानांत
$A \lif B$ या रूपाची
सूत्रे असतात.

नैसर्गिक निगमनाच्या
पहिल्या दोन आवृत्त्या
१९३० च्या दशकात गरहार्ड
गेंटझेन आणि स्टॅनिस्लॉ
यास्कोव्हस्की यांनी मांडल्या.
सिद्धता-उपपत्तीचे अभ्यासक
बहुधा गेंट्सेनच्या पद्धतीची
चर्चा करतात; अध्यापनात
सर्वाधिक वापरल्या जाणाऱ्या
पद्धतींचा उगम याश्कोव्स्कीच्या
पद्धतीत आहे. दोन्हीत
\emph{गृहीतकांपासून} सुरुवात
करता येते: अशी सूत्रे
की ज्यांना सिद्ध करावे
लागत नाही, पण इतर
सूत्रे सिद्ध करणे
करताना वापरता येते.
संपूर्ण सिद्धता अशा
गृहीतकांवर अवलंबून असू शकते.
ही गृहीतके असिद्ध असल्यामुळे
ती सिद्धता तिचा निष्कर्ष
(तिचे \emph{अंतिम-सूत्र})
सिद्ध करत नाही; निष्कर्ष
गृहीतकांपासून निष्पन्न होतो,
एवढेच ती दाखवते.

काही अनुमान नियमांचा
निष्कर्ष काही गृहीतकांवर
आता अवलंबून राहत नाही:
असे नियम गृहीतके “मुक्त”
करतात. उदाहरणार्थ,
\Intro{\lif} नियम
गृहीतक~$A$ पासून
निष्कर्ष~$B$ ची
सिद्धता घेऊन
$A \lif B$ हा
निष्कर्ष देतो.
$A \lif B$ मध्ये
संपणारी सिद्धता
आता~$A$ वर अवलंबून
नसते: ते गृहीतक
मुक्त झालेले असते.
गेंट्सेनच्या पद्धतीत
सिद्धता हे
सूत्रांचे वृक्ष
असतात; गृहीतके सर्वात
वरची सूत्रे असतात,
आणि \Intro{\lif}
सारख्या नियमांवरील खुणा
कोणती गृहीतके मुक्त होतात
ते दाखवतात.
याश्कोव्स्कीच्या पद्धतीत
एखाद्या गृहीतकावर अवलंबून
असलेले सिद्धता चे भाग
वेगळे दाखवतात: उदाहरणार्थ,
उभी रेषा लावून आत सरकवलेले
किंवा चौकटीत बंद केलेले.
त्या चौकटीची पहिली ओळ
गृहीतक~$A$ असते;
अंतिम ओळ सोपाधिक विधानाचे
उत्तरांग $B$ असते;
$A$ मुक्त केले जाते;
आणि \Intro{\lif} चा
निष्कर्ष~$A \lif B$
चौकटीबाहेर असतो.

या पद्धती “नैसर्गिक”
आहेत, कारण त्यांचे अनुमान
नियम आपण (निदान गणितज्ञ)
सहजपणे जसे युक्तिवाद करतो
त्याशी जुळतात. एखादे
सोपाधिक विधान सिद्ध
करण्यासाठी त्याचे पूर्वांग
गृहीत धरून त्यावरून उत्तरांग
सिद्ध करतो: तो \Intro{\lif}.
सोपाधिक विधान $A \lif B$
आणि त्याचे पूर्वांग~$A$
माहीत असेल, तर~$B$
निष्पन्न करतो: तो
\Elim{\lif}.
विकल्पयोग~$A \lor B$
गृहीत धरून काही निष्कर्ष
दाखवण्यासाठी प्रकरणांनुसार
सिद्धता वापरतो: तो
\Elim{\lor}.
सार्वत्रिक विधान
$\lforall[x][A(x)]$
सिद्ध करण्यासाठी
कोणताही~$c$ घेऊन
$A(c)$ सिद्ध करतो:
तो \Intro{\lforall}.
अशाच प्रकारे पुढे जाता येते.

उदाहरणार्थ, गेंट्सेन-पद्धतीच्या
नैसर्गिक निगमनात
$A \lif (A \lor B)$
ची ही सोपी सिद्धता पाहा:
\[
\AxiomC{$\Discharge{A}{x}$}
\RightLabel{\Intro{\lor}}
\UnaryInfC{$A \lor B$}
\DischargeRule{\Intro{\lif}}{x}
\UnaryInfC{$A \lif (A \lor B)$}
\DisplayProof
\]
ती गृहीतक~$A$ पासून
सुरू होते; \Intro{\lor}
वापरून~$A \lor B$
निष्पन्न करते; मग
\Intro{\lif} वापरून
$A \lif (A \lor B)$
निष्पन्न करते.
\Intro{\lif} अनुमान
गृहीतक~$A$ मुक्त करते;
खूण~$x$ ते दर्शवते.

सिद्धता-उपपत्तीचे अभ्यासक
सूत्रांच्या वृक्षांवर
काम करणाऱ्या गेंट्सेनच्या
मूळ पद्धती पसंत करतात;
तरी एक रूपांतरही
सिद्धता-उपपत्तीत महत्त्वाचे आहे.
गृहीतके~$\Gamma$ यांपासून
$A$ ची सिद्धता
म्हणजे $A$ हे~$\Gamma$
पासून निष्पन्न होते,
अर्थात $\Gamma \Entails A$,
हे दाखवते.
नैसर्गिक निगमनाचे अनुमान
नियम अशा निष्पन्नतेबद्दल
विधान करतात असेही
स्वाभाविकपणे वाचता येते.
उदाहरणार्थ, \Intro{\lif}
म्हणतो की जर $\Gamma + A
\Entails B$, तर
$\Gamma \Entails A \lif B$.
म्हणून नियम थेट गृहीतके
आणि त्यांपासून निष्पन्न
होणारा निष्कर्ष या दोहोंवर
लागू होतील असे त्यांचे
रूप देता येते: म्हणजे
ते $\Gamma \Sequent A$
या क्रमवर्तींवर काम करतात.
वरील सोपी सिद्धता अशा
पद्धतीत पुढीलप्रमाणे दिसेल:
\[
\Axiom$x:A \fCenter A$
\RightLabel{\Intro{\lor}}
\UnaryInf$x:A \fCenter A \lor B$
\DischargeRule{\Intro{\lif}}{x}
\UnaryInf$\fCenter A \lif (A \lor B)$
\DisplayProof
\]
येथे नियम तेच आहेत:
ते~$\Sequent$ च्या
उजवीकडील सूत्र
वर काम करतात;
डावीकडील सूत्रे
प्रत्येक टप्प्यावर कोणत्या
गृहीतकांवर अवलंबित्व आहे
ते फक्त दाखवतात.

परिचय/निरसन
या नियम-जोड्यांचा पूर्ण संच
स्वतःहून अभिजात तर्कशास्त्रासाठी
पूर्ण नसतो. त्यात~$\lfalse$
शी संबंधित आणखी दोन
नियम जोडावे लागतात.
पहिला “अंतःप्रज्ञावादी
असत्यता नियम”~\FalseInt,
म्हणजे “स्फोट”:
$\lfalse$ पासून काहीही
निष्पन्न करता येते.
दुसरा अप्रत्यक्ष सिद्धतेचा
अभिजात सिद्धांत~\FalseCl{}
(किंवा “अभिजात असत्यता नियम”):
$\lnot A$ पासून
$\lfalse$ सिद्ध
करता येत असेल, तर
$A$ सिद्ध झाले.
\FalseCl वगळल्यास
अंतःप्रज्ञावादी तर्कशास्त्राला
योग्य अशी पद्धत मिळते.
दोन्ही नियम वगळल्यास
“किमान तर्कशास्त्र”
मिळते.

अभिजात आणि
अंतःप्रज्ञावादी तर्कशास्त्रासाठी
गेंट्सेन-पद्धतीच्या
नैसर्गिक निगमनाची आपण
चर्चा करू. या
\Log{N1c} आणि \Log{N1i}
पद्धती आहेत.
\Log{N2c} आणि \Log{N2i}
या संबंधित पद्धती
केवळ सूत्रे~$A$
यांवर काम करण्याऐवजी
$\Gamma \Sequent A$
या क्रमवर्तींवर काम करतात.












\section{नियम आणि सिद्धता}\label{pt:nat:rp:sec}

प्रथम काही व्याख्या करू
आणि संज्ञा निश्चित करू.

~\Log{N1c} आणि~\Log{N1i}
मधील $A \lif B$
उल्लेखणारे दोन नियम पाहू:
\[
\bottomAlignProof
\AxiomC{$\Discharge{A}{x}$}
\shortDeduce
\DeduceC{$B$}
\DischargeRule{\Intro{\lif}}{x}
\UnaryInfC{$A \lif B$}
\DisplayProof
\qquad
\bottomAlignProof
\AxiomC{$A \lif B$}
\AxiomC{$A$}
\RightLabel{\Elim{\lif}}
\BinaryInfC{$B$}
\DisplayProof
\]
\Elim{\lif} नियम सोपा आहे.
रेषेच्या वरची सूत्रे
ही \emph{आधारविधाने} होत.
डावीकडील सूत्र
हे \emph{मुख्य}
सूत्र~$A \lif B$
आहे. सर्व निरसन
नियमांत असे एक आधारविधान
असते आणि ते नेहमी
सर्वात डावीकडे असते.
त्याला अनुमानाचे
\emph{प्रमुख} आधारविधान
म्हणतात. या उदाहरणाप्रमाणे
इतर आधारविधाने असतील,
तर त्यांना \emph{गौण}
आधारविधाने म्हणतात.
परिचय नियमांची
आधारविधानेही प्रमुख
आधारविधाने म्हणतात.
रेषेखालील सूत्र
हा \emph{निष्कर्ष} आहे.

\Intro{\lif} नियमाला
एकच आधारविधान असते;
येथे निष्कर्ष हे
मुख्य सूत्र~$A \lif B$
असते. सर्व परिचय
नियम असेच असतात:
निष्कर्ष हे मुख्य सूत्र
असते, आणि त्याचा
मुख्य संकारकच “प्रवेशित”
केला जातो.

~$\Discharge{A}{x}$
या लेखनाचा अर्थ पुढीलप्रमाणे.
एखाद्या सिद्धता मध्ये
\Intro{\lif} अनुमान
आधारविधान म्हणून
सूत्र~$B$
वर लावतात.
~$B$ मध्ये संपणाऱ्या
उप-सिद्धता च्या
सुरुवातीला काही
सूत्रे असतात;
त्यांना \emph{गृहीतके}
म्हणतात. ~$A$
रूपाची एक किंवा अधिक
गृहीतके धरून $B$
ची सिद्धता म्हणजे
$A$ वरून~$B$
निष्पन्न होते हे दाखवणारी
सिद्धता होय.
\Intro{\lif} लावल्यावर
$A \lif B$
निष्पन्न होते; आता
निष्कर्ष गृहीतक~$A$
वर अवलंबून राहत नाही.
हे दाखवण्यासाठी,
\Intro{\lif} असलेल्या
सिद्धता मधील
$A$ रूपाची संबंधित
गृहीतके \emph{मुक्त}
किंवा \emph{बंद} केली
जातात. कोणती गृहीतके
मुक्त करता येतात, हे
नियम सांगतो.
$A \lif B$ निष्कर्ष
असलेल्या \Intro{\lif}
मध्ये ती $A$ रूपाची,
म्हणजे सोपाधिक विधानाच्या
पूर्वांगाशी जुळणारी
गृहीतके असतात.
कोणत्या नियमाने कोणती
गृहीतके मुक्त होतात हे
दाखवणे कधी महत्त्वाचे
असते; त्यासाठी मुक्तीची
खूण~$x$ वापरतात.
मात्र गृहीतके मुक्त
करणाऱ्या नियमाने योग्य
रूपाची \emph{सर्वच}
गृहीतके, किंवा एकही
गृहीतक, मुक्त करणे
\emph{बंधनकारक नाही}.
उदाहरणार्थ, पहिल्या
\Intro{\lif} ने एकही
गृहीतक मुक्त केलेले
नसले तरी पुढील सिद्धता
योग्य आहे:
\[
\AxiomC{$\Discharge{A}{y}$}
\DischargeRule{\Intro{\lif}}{x}
\UnaryInfC{$B \lif A$}
\DischargeRule{\Intro{\lif}}{y}
\UnaryInfC{$A \lif (B \lif A)$}
\DisplayProof
\]
नियमाने एकही गृहीतक
मुक्त केले नाही, तर
मुक्तीची खूण वगळतो
(येथेही तसे केले आहे).

पुढील सिद्धता पाहा:
\[
\AxiomC{$\Discharge{A}{x}$}
\AxiomC{$\Discharge{A}{y}$}
\RightLabel{\Intro{\land}}
\BinaryInfC{$A \land A$}
\RightLabel{\Elim{\land}}
\UnaryInfC{$A$}
\DischargeRule{\Intro{\lif}}{x}
\UnaryInfC{$A \lif A$}
\DischargeRule{\Intro{\lif}}{y}
\UnaryInfC{$A \lif (A \lif A)$}
\DisplayProof
\]
पहिले \Intro{\lif}
अनुमान डावीकडील
गृहीतक~$A^x$ मुक्त करते;
दुसरे उजवीकडील
गृहीतक~$A^y$ मुक्त करते.
म्हणजे खूण~$x$ असलेली
सर्व गृहीतके खूण~$x$
असलेल्या \Intro{\lif}
ने, आणि खूण~$y$
असलेली सर्व गृहीतके
खूण~$y$ असलेल्या
त्या अनुमानाने
मुक्त होतात.
यासाठी एकाच खूणेची
सर्व गृहीतके तेच
सूत्र असणे
महत्त्वाचे आहे.

संख्यापकांचे नियम
किंचित अधिक गुंतागुंतीचे आहेत.
उदाहरणार्थ,~\Log{N1i}
मध्ये $\lexists$
साठीचे नियम असे:
\[
\bottomAlignProof
\AxiomC{$A(t)$}
\RightLabel{\Intro{\lexists}}
\UnaryInfC{$\lexists[x][A(x)]$}
\DisplayProof
\qquad
\bottomAlignProof
\AxiomC{$\lexists[x][A(x)]$}
\AxiomC{$\Discharge{A(c)}{x}$}
\shortDeduce
\DeduceC{$B$}
\DischargeRule{\Elim{\lexists}}{x}
\BinaryInfC{$B$} \DisplayProof
\bottomAlignProof
\]
\Intro{\lexists} मध्ये
आधारविधान~$A(t)$ असते;
येथे $t$ हे बंद पद,
म्हणजे चल नसलेले पद आहे.
असे अनुमान योग्य,
म्हणजे \Intro{\lexists}
या नियम-रूपबंधाचे उदाहरण,
तेव्हाच असते जेव्हा
एखादे सूत्र~$A$,
चल~$x$ आणि बंद पद~$t$
असते ज्यामुळे निष्कर्ष
$\lexists[x][A]$
आणि आधारविधान
~$\Subst{A}{t}{x}$
असते. \Elim{\lexists}
नियमाने $A(c)$
रूपाची गृहीतके मुक्त
करता येतात: प्रत्येक
अनुमानासाठी असा एखादा
स्थिरांक~$c$ असतो
की $\Subst{A}{c}{x}$
रूपाची गृहीतके मुक्त
होतात. एका अनुमानाने
मुक्त होणाऱ्या सर्व
गृहीतकांत तोच~$c$
वापरला पाहिजे.
या स्थिरांकाला
\Elim{\lexists}
अनुमानाचे
\emph{आयगेन चर} म्हणतात.
~$c$ हे मुक्त होणाऱ्या
~$A(c)$ व्यतिरिक्त
इतर उघड्या गृहीतकात,
~$B$ मध्ये किंवा
~$A$ मध्ये येत
नसेल तेव्हाच अनुमान
योग्य असते. याला
\emph{आयगेन चराची अट}
म्हणतात.

~\Log{N1c} आणि~\Log{N1i}
मधील सिद्धता म्हणजे
नियम वापरून स्वयंसिद्धकांपासून
विगमनाने तयार होणारे
सूत्रांचे सांत वृक्ष.
प्रत्येक पान हे मुक्तीच्या
खुणेसह किंवा तिच्याविना
असलेले गृहीतक असते;
इतर प्रत्येक सूत्र
त्याच्या लगत वरच्या
सूत्रांपासून
पद्धतीच्या अनुमान नियमाने
निष्पन्न होते.
अनुमानाच्या वरच्या
वृक्षातील गृहीतक म्हणून
आलेले एखादे सूत्र
त्या अनुमानापर्यंत मुक्त
झालेले नसेल, तर त्याला
(त्या अनुमानाशी संबंधित)
\emph{उघडे} म्हणतात.

\begin{defn}\label{pt:nat:rp:defn:proof}
~\Log{N1c} आणि~\Log{N1i}
मधील \emph{सिद्धता}
हे पुढीलप्रमाणे विगमनाने
परिभाषित सूत्रे
चे सांत वृक्ष आहेत:
\begin{enumerate}
  \item\label{pt:nat:rp:defn:prf:assumption}
  मुक्तीची $x$ किंवा $y$
  सारखी खूण असलेले
  प्रत्येक सूत्र
  स्वतःच सिद्धता असते.
  एकाच सूत्र
  ची सिद्धता असेल,
  तर ते सूत्र
  उघडे गृहीतक मानतात.
  \item\label{pt:nat:rp:defn:prf:one-prem}
  $R$ हा एक, दोन किंवा
  तीन आधारविधाने
  $A_1$, $A_2$, $A_3$
  आणि निष्कर्ष~$A$
  असलेला नियम असू द्या.
  $\delta_1$, $\delta_2$,
  $\delta_3$ या अनुक्रमे
  $A_1$, $A_2$,
  $A_3$ मध्ये संपणाऱ्या
  सिद्धता असतील, तर
  पुढीलपैकी योग्य रचना
  अनुक्रमे सिद्धता असते:
  \[
  \AxiomC{}
  \RightLabel{$\delta_1$}
  \DeduceC{$A_1$}
  \RightLabel{$R$}
  \UnaryInfC{$A$}
  \DisplayProof
\qquad
  \AxiomC{}
  \RightLabel{$\delta_1$}
  \DeduceC{$A_1$}
  \AxiomC{}
  \RightLabel{$\delta_2$}
  \DeduceC{$A_2$}
  \RightLabel{$R$}
  \BinaryInfC{$A$}
  \DisplayProof
\qquad
  \AxiomC{}
  \RightLabel{$\delta_1$}
  \DeduceC{$A_1$}
  \AxiomC{}
  \RightLabel{$\delta_2$}
  \DeduceC{$A_2$}
  \AxiomC{}
  \RightLabel{$\delta_3$}
  \DeduceC{$A_3$}
  \RightLabel{$R$}
  \TrinaryInfC{$A$}
  \DisplayProof
  \]
  मात्र त्या सिद्धता
  मधील गृहीतक-खुणा
  सुसंगत असल्या पाहिजेत
  (दोन भिन्न
  सूत्रे ना
  तीच मुक्तीची खूण
  नसावी), आणि ही
  अनुमाने~$R$ चे
  योग्य उपयोग असले पाहिजेत.
  \item नव्या सिद्धता
  मधील सर्वात वरची
  सूत्र उदाहरणे
  ही त्या सिद्धता
  ची \emph{उघडी गृहीतके}
  मानली जातात, जर ती
  $\delta_1$, $\delta_2$
  किंवा~$\delta_3$
  मधील उघडी गृहीतके
  असतील आणि नियमाने
  मुक्त केलेली नसतील.
  नियमावर खूण~$x$
  (किंवा $x\, y$)
  असेल, तर \Intro{\lif}
  आणि~\Intro{\lnot}
  साठी $\delta_1$
  मधील खूण~$x$
  असलेली सर्व उघडी
  गृहीतके मुक्त होतात;
  \Elim{\lexists}
  साठी $\delta_2$
  मधील खूण~$x$
  असलेली सर्व उघडी
  गृहीतके मुक्त होतात;
  \Elim{\lor} साठी
  $\delta_2$ मधील
  खूण~$x$ असलेली
  आणि~$\delta_3$ मधील
  खूण~$y$ असलेली
  सर्व उघडी गृहीतके
  मुक्त होतात.
\end{enumerate}
\end{defn}

सिद्धता ची
विगमनात्मक रचना करताना
वापरलेल्या सिद्धता ना
तिच्या \emph{उप-सिद्धता}
म्हणतात. अनुमानाच्या
आधारविधानांत संपणाऱ्या
उप-सिद्धता ना त्या
अनुमानाच्या \emph{निकटच्या
उप-सिद्धता} म्हणतात.
\Log{N1c} आणि \Log{N1i}
यांचे नियम
\ref{pt:nat:rp:tab:N1} मध्ये आहेत.








\begin{table}
\begin{center}
\begin{tabular}[t]{@{}cc@{}}
\bottomAlignProof
\AxiomC{$\Discharge{A}{x}$}
\shortDeduce
\DeduceC{$B$}
\DischargeRule{\Intro{\lif}}{x}
\UnaryInfC{$A \lif B$}
\DisplayProof
&
\bottomAlignProof
\AxiomC{$A \lif B$}
\AxiomC{$A$}
\RightLabel{\Elim{\lif}}
\BinaryInfC{$B$}
\DisplayProof
\\[4ex]
\bottomAlignProof
\AxiomC{$A$}
\AxiomC{$B$}
\RightLabel{\Intro{\land}}
\BinaryInfC{$A \land B$}
\DisplayProof
&
\begin{tabular}[b]{@{}l@{}}
\AxiomC{$A \land B$}
\RightLabel{\Elim{\land}[1]}
\UnaryInfC{$A$}
\DisplayProof
\\[3ex]
\bottomAlignProof
\AxiomC{$A \land B$}
\RightLabel{\Elim{\land}[2]}
\UnaryInfC{$B$}
\DisplayProof
\end{tabular}
\\[4ex]
\begin{tabular}{@{}l@{}}
\AxiomC{$A$}
\RightLabel{\Intro{\lor}[1]}
\UnaryInfC{$A \lor B$}
\DisplayProof
\\[3ex]
\AxiomC{$B$}
\RightLabel{\Intro{\lor}[2]}
\UnaryInfC{$A \lor B$}
\DisplayProof
\end{tabular}
&
\AxiomC{$A \lor B$}
\AxiomC{$\Discharge{A}{x}$}
\shortDeduce
\DeduceC{$C$}
\AxiomC{$\Discharge{B}{y}$}
\shortDeduce
\DeduceC{$C$}
\DischargeRule{\Elim{\lor}}{x\,y}
\TrinaryInfC{$C$}
\DisplayProof
\bottomAlignProof
\\[6ex]
\bottomAlignProof
\AxiomC{$\Discharge{A}{x}$}
\shortDeduce
\DeduceC{$\lfalse$}
\DischargeRule{\Intro{\lnot}}{x}
\UnaryInfC{$\lnot A$}
\DisplayProof
&
\bottomAlignProof
\AxiomC{$\lnot A$}
\AxiomC{$A$}
\RightLabel{\Elim{\lnot}}
\BinaryInfC{$\lfalse$}
\DisplayProof
\\[4ex]
\bottomAlignProof
\AxiomC{$A(a)$}
\RightLabel{\Intro{\lforall}}
\UnaryInfC{$\lforall[x][A(x)]$}
\DisplayProof
&
\bottomAlignProof
\AxiomC{$\lforall[x][A(x)]$}
\RightLabel{\Elim{\lforall}}
\UnaryInfC{$A(t)$}
\DisplayProof
\\[4ex]
\bottomAlignProof
\AxiomC{$A(t)$}
\RightLabel{\Intro{\lexists}}
\UnaryInfC{$\lexists[x][A(x)]$}
\DisplayProof
&
\bottomAlignProof
\AxiomC{$\lexists[x][A(x)]$}
\AxiomC{$\Discharge{A(a)}{x}$}
\shortDeduce
\DeduceC{$C$}
\DischargeRule{\Elim{\lexists}}{x}
\BinaryInfC{$C$}
\DisplayProof
\\[4ex]
\bottomAlignProof
\AxiomC{$\lfalse$}
\RightLabel{\FalseInt}
\UnaryInfC{$A$}
\DisplayProof
&
\bottomAlignProof
\AxiomC{$\Discharge{\lnot A}{x}$}
\shortDeduce
\DeduceC{$\lfalse$}
\DischargeRule{\FalseCl}{x}
\UnaryInfC{$A$}
\DisplayProof
\end{tabular}
\end{center}
\caption{\Log{N1i} आणि \Log{N1c}
यांचे नियम. $\Intro{\lforall}$
मध्ये स्थिरांक~$a$
निष्कर्षात किंवा उघड्या
गृहीतकात येता कामा नये.
$\Elim{\lexists}$ मध्ये
$a$ दोन्ही आधारविधानांच्या
निष्कर्षांत किंवा अनुमानानंतर
उरलेल्या उघड्या गृहीतकांत
येता कामा नये.
\Log{N1i} मध्ये \FalseCl
नसतो.}
\label{pt:nat:rp:tab:N1}
\end{table}




\begin{defn}
सिद्धता~$\delta$ ची
\emph{उंची}~
$\pheight{\delta}$
पुढीलप्रमाणे विगमनाने
परिभाषित करतात.
\begin{enumerate}
  \item $\delta$ केवळ
  गृहीतकाची असेल, तर
  $\pheight{\delta} = 0$.
  \item $\delta$ चा शेवट
  एका आधारविधानाच्या
  अनुमानात होत असेल
  आणि निकटची
  उप-सिद्धता~$\delta_1$
  असेल, तर
  $\pheight{\delta} =
  \pheight{\delta_1} + 1$.
  \item $\delta$ चा शेवट
  दोन आधारविधानांच्या
  अनुमानात होत असेल
  आणि निकटच्या
  उप-सिद्धता~$\delta_1$
  आणि~$\delta_2$ असतील,
  तर $\pheight{\delta} =
  \max(\pheight{\delta_1},
  \pheight{\delta_2}) + 1$.
  \item $\delta$ चा शेवट
  \Elim{\lor} मध्ये होत
  असेल आणि निकटच्या
  उप-सिद्धता~$\delta_1$,
  $\delta_2$ आणि~$\delta_3$
  असतील, तर
  $\pheight{\delta} =
  \max(\pheight{\delta_1},
  \pheight{\delta_2},
  \pheight{\delta_3}) + 1$.
\end{enumerate}
क्रमवर्ती~$S$ ची उंची
$\le n$ असलेली
सिद्धता असेल,
तर $ \Proves[n] S$
असे लिहितो.
\end{defn}

\begin{ex}
~\Log{N1i} मध्ये कोणतेही
सूत्र हे स्वतःपासून
उघडे गृहीतक घेऊन
सिद्धता असते.
म्हणून
\[
A \land B^x
\]
ही सिद्धता~$\delta_1$
आहे; तिची उंची~$0$.
~$\delta_1$ वर
~\Elim{\land} च्या
दोनपैकी एक आवृत्ती
लावून सिद्धता~$\delta_2$
आणि~$\delta_3$ मिळतात:
\[
\AxiomC{$A \land B^x$}
\RightLabel{\Elim{\land}[2]}
\UnaryInfC{$B$}
\DisplayProof
\qquad
\AxiomC{$A \land B^x$}
\RightLabel{\Elim{\land}[1]}
\UnaryInfC{$A$}
\DisplayProof
\]
~$\delta_2$ आणि
~$\delta_3$ मधील
अंतिम अनुमानांच्या निकटच्या
उप-सिद्धता सारख्याच:
$A \land B$ या
गृहीतकापासून तयार झालेल्या.
दोन्ही सिद्धता मध्ये
$A\land B$ हे
उघडे गृहीतक म्हणून येते.
आपल्याला
$\pheight{\delta_2} =
\pheight{\delta_3} = 1$
मिळते.

~\Intro{\land} नियमात
$B$ आणि $A$
या रूपाची दोन
आधारविधाने असतील,
तर $B \land A$
या रूपाचा निष्कर्ष
निघतो. त्यामुळे
पुढील सिद्धता~$\delta_4$
आहे:
\[
\AxiomC{$A \land B^x$}
\RightLabel{\Elim{\land}[2]}
\UnaryInfC{$B$}
\LeftSubproofLabel{$\delta_2$}
\AxiomC{$A \land B^x$}
\RightLabel{\Elim{\land}[1]}
\UnaryInfC{$A$}
\RightSubproofLabel{$\delta_3$}
\RightLabel{\Intro{\land}}
\BinaryInfC{$B \land A$}
\DisplayProof
\]
तिची उंची
$\pheight{\delta_4} =
\max(\pheight{\delta_2},
\pheight{\delta_3})+1
= \max(1,1) + 1 = 2$.
तिच्यात $A \land B$
या गृहीतकाची दोन
उदाहरणे आहेत; दोन्ही
उघडी आहेत.

शेवटी~\Intro{\lif}
लावून $\delta_5$
मिळवू:
\[
\AxiomC{$\Discharge{A \land B}{x}$}
\RightLabel{\Elim{\land}[2]}
\UnaryInfC{$B$}
\LeftSubproofLabel{$\delta_2$}
\AxiomC{$\Discharge{A \land B}{x}$}
\RightLabel{\Elim{\land}[1]}
\UnaryInfC{$A$}
\RightSubproofLabel{$\delta_3$}
\RightLabel{\Intro{\land}}
\BinaryInfC{$B \land A$}
\DischargeRule{\Intro{\lif}}{x}
\UnaryInfC{$(A \land B) \lif (B \land A)$}
\DisplayProof
\]
तिची उंची
$\pheight{\delta_4} + 1 = 3$.
\Intro{\lif} अनुमान
$A \land B$
या रूपाची, खूण~$x$
असलेली सर्व गृहीतके
मुक्त करते; म्हणजे
निकटच्या उपसिद्धतेतील
आधी उघडी असलेली
दोन्ही गृहीतके.
~$\delta_5$ ची
तीच दोन गृहीतके असल्याने
तिच्यात उघडे गृहीतक
उरत नाही; म्हणून ती
तिच्या अंतिम-सूत्र
$(A \land B)
\lif (B \land A)$
\emph{ची} सिद्धता
आहे.
\end{ex}












\section{क्रमवर्तींसह नैसर्गिक निगमन: \Log{N2c} आणि~\Log{N2i}}\label{pt:nat:seq:sec}

~\Log{N1c} किंवा
\Log{N1i} मधील
सिद्धता~$\delta$
तिचे अंतिम-सूत्र~$A$
हे तिच्या उघड्या गृहीतकांपासून
निष्पन्न होते हे दाखवते.
~$\delta$ मध्ये $B^x$
हे उघडे गृहीतक असलेल्या
सूत्रे~$B$
यांचा संच~$\Gamma$
असेल, तर हे
$\delta \Proves \Gamma
\Sequent A$ असे
लिहिता येते.
यावरून क्रमवर्तींसाठीही
नैसर्गिक निगमन पद्धत
बांधता येईल असे दिसते.
अशा पद्धतीत प्रत्येक
क्रमवर्ती तिच्या उजवीकडील
सूत्र कोणत्या
उघड्या गृहीतकांवर अवलंबून
आहे, म्हणजे त्या
क्रमवर्तीत संपणाऱ्या
उप-सिद्धता मध्ये
कोणती गृहीतके उघडी
आहेत, हे सांगते.
यालाच क्रमवर्तींसह,
किंवा “क्रमवर्ती-पद्धतीचे”
नैसर्गिक निगमन म्हणतो.
मिळणाऱ्या पद्धती
\Log{N2c} आणि~\Log{N2i}
अशा आहेत.

~\Log{N1} आणि
\Log{N2} यांतील
संबंध अधिक निकट करण्यासाठी
डावीकडील संच~$\Gamma$
हे खूण लावलेल्या
सूत्रे~$x:B$
यांचे धरू. ~\Log{N1}
चे नियम केवळ
आधारविधानांतील
सूत्रे वर,
म्हणजे निकटच्या
उप-सिद्धता च्या
अंतिम-सूत्रे वर
लागतात. त्यामुळे
\Log{N2} चे नियम
क्रमवर्तीच्या उजव्या
भागावरच लागतात.
क्रमवर्तींच्या डावीकडील
खूण लावलेल्या सूत्रांना
म्हणून फक्त संदर्भ
म्हणता येते.

~\Log{N2} सिद्धतांत
गृहीतकांच्या जागी
सर्वात वर अशा रूपाची
स्वयंसिद्धक-क्रमवर्ती असते:
\[x:A \Sequent A.\]
नियम~\Log{N1}
च्या नियमांशी तंतोतंत
जुळतात; ते
\ref{pt:nat:seq:tab:N2} मध्ये आहेत.








\begin{table}
\begin{tabular}{@{}c@{}c@{}}
\Axiom$x: A, \Gamma \fCenter B$
\RightLabel{\Intro{\lif}}
\UnaryInf$\Gamma \fCenter A \lif B $
\DisplayProof
&
\Axiom$\Gamma_1 \fCenter A \lif B$
\Axiom$\Gamma_2 \fCenter A$
\RightLabel{\Elim{\lif}}
\BinaryInf$\Gamma_1, \Gamma_2 \fCenter B$
\DisplayProof
\\[4ex]
\Axiom$\Gamma_1 \fCenter A$
\Axiom$\Gamma_2 \fCenter B$
\RightLabel{\Intro{\land}}
\BinaryInf$\Gamma_1, \Gamma_2 \fCenter A \land B $
\DisplayProof
&
\begin{tabular}[b]{@{}l@{}}
\Axiom$\Gamma \fCenter A \land B$
\RightLabel{\Elim{\land}[1]}
\UnaryInf$\Gamma \fCenter A$
\DisplayProof
\\[3ex]
\Axiom$\Gamma \fCenter A \land B$
\RightLabel{\Elim{\land}[2]}
\UnaryInf$\Gamma \fCenter B$
\DisplayProof\\[3ex]
\end{tabular}
\\[4ex]
\Axiom$\Gamma \fCenter A$
\RightLabel{\Intro{\lor}[1]}
\UnaryInf$\Gamma \fCenter A \lor B$
\DisplayProof
&
\Axiom$ \Gamma \fCenter B$
\RightLabel{\Intro{\lor}[2]}
\UnaryInf$\Gamma \fCenter A \lor B$
\DisplayProof\\[3ex]
\multicolumn{2}{c}{
\begin{tabular}[b]{@{}l@{}}
\Axiom$\Gamma_1 \fCenter A \lor B$
\Axiom$x:A, \Gamma_2 \fCenter C$
\Axiom$y:B, \Gamma_3 \fCenter C$
\RightLabel{\Elim{\lor}}
\TrinaryInf$\Gamma_1, \Gamma_2, \Gamma_3 \fCenter C$
\DisplayProof
\end{tabular}
}
\\[4ex]
\Axiom$x:A, \Gamma \fCenter \lfalse$
\RightLabel{\Intro{\lnot}}
\UnaryInf$ \Gamma \fCenter \lnot A $
\DisplayProof
&
\Axiom$\Gamma_1 \fCenter \lnot A $
\Axiom$\Gamma_2 \fCenter A $
\RightLabel{\Elim{\lnot}}
\BinaryInf$\Gamma_1, \Gamma_2 \fCenter \lfalse$
\DisplayProof
\\[4ex]
\Axiom$\Gamma \fCenter A(a) $
\RightLabel{\Intro{\lforall}}
\UnaryInf$\Gamma \fCenter \lforall[x][A(x)]$
\DisplayProof
&
\Axiom$\Gamma \fCenter \lforall[x][A(x)]$
\RightLabel{\Elim{\lforall}}
\UnaryInf$\Gamma \fCenter A(t)$
\DisplayProof
\\[4ex]
\Axiom$\Gamma \fCenter A(t) $
\RightLabel{\Intro{\lexists}}
\UnaryInf$\Gamma \fCenter \lexists[x][A(x)]$
\DisplayProof
&
\Axiom$\Gamma_1 \fCenter  \lexists[x][A(x)]$
\Axiom$x: A(a), \Gamma_2 \fCenter C$
\RightLabel{\Elim{\lexists}}
\BinaryInf$\Gamma_1, \Gamma_2 \fCenter C$
\DisplayProof
\\[4ex]
\Axiom$\Gamma \fCenter \lfalse$
\RightLabel{\FalseInt}
\UnaryInf$\Gamma \fCenter A$
\DisplayProof
&
\Axiom$x: \lnot A, \Gamma \fCenter \lfalse$
\RightLabel{\FalseCl}
\UnaryInf$\Gamma \fCenter A$
\DisplayProof
\end{tabular}
\caption{\Log{N2c} आणि \Log{N2i}
यांचे नियम. $\Intro{\forall}$
आणि $\Elim{\exists}$ मध्ये
स्थिरांक~$a$
निष्कर्ष-क्रमवर्तीमध्ये
येता कामा नये.
$\Gamma_i$ हे खूण
लावलेल्या सूत्रे~$x :
A$ चे संच आहेत.
\Intro{\lforall} आणि
\Elim{\lexists} नियमांसाठी
आयगेन चराची अट पूर्ण
झाली पाहिजे:
$a$ निष्कर्ष-क्रमवर्तीमध्ये
येता कामा नये.
\Log{N2i} मध्ये \FalseCl
नसतो.}
\label{pt:nat:seq:tab:N2}
\end{table}




~\Log{N1c} आणि
\Log{N1i} मध्ये
\Intro{\lif},
\Intro{\lnot},
\Elim{\lor} आणि
\Elim{\lexists} हे
नियम गृहीतके मुक्त करू
शकतात, पण तसे करणे
बंधनकारक नाही.
म्हणून~\Log{N2c}
आणि~\Log{N2i}
मध्येही संबंधित आधारविधानाच्या
संदर्भात अपेक्षित
सूत्रे $A$,
$B$ आणि $A(a)$
नसली तरी त्या
नियमांचा योग्य उपयोग
मान्य करतो.

\begin{prop}
$\delta$ ही~\Log{N1c}
किंवा \Log{N1i}
मध्ये~$A$ ची
सिद्धता असेल,
तर~\Log{N2c}
(\Log{N2i}) मध्ये
$\Gamma \Sequent A$
ची सिद्धता~$\delta'$
असते. येथे~$\Gamma$
हा असा सर्व~$x:B$
यांचा संच आहे की
$B^x$ हे~$\delta$
चे उघडे गृहीतक असते.
\end{prop}

\begin{proof}
  $n=\pheight{\delta}$
  वर विगमन करू.
  $n=0$ असेल, तर
  $\delta$ ही एकच
  उघडे गृहीतक~$A^x$.
  मग $\Gamma = \{x:A\}$
  आणि $\Gamma
  \Sequent A$ हे
  \Log{N2c} किंवा
  ~\Log{N2i} चे
  स्वयंसिद्धक आहे.

  $n>0$ असेल, तर
  ~$\delta$ मधील अंतिम
  अनुमानानुसार प्रसंग
  पाडू. येथे एकच प्रसंग
  पाहू; उरलेले सरावासाठी.

  अंतिम अनुमान
  $\Intro{\lif}$ असल्याचे
  समजा. मग~$\delta$
  चा शेवट असा होतो:
  \[
  \AxiomC{$\Discharge{B}{x}$}
  \RightLabel{$\delta_1$}
  \DeduceC{$C$}
  \DischargeRule{\Intro{\lif}}{x}
  \UnaryInfC{$B \lif C$}
  \DisplayProof
  \]
  विगमन गृहीतकानुसार
  ~\Log{N2c} (\Log{N2i})
  मध्ये
  $\Gamma' \Sequent C$
  ची सिद्धता~$\delta_1'$
  आहे; $\Gamma'$ म्हणजे
  ~$\delta_1$ मधील
  खूण लावलेली उघडी
  गृहीतके. $B^x$
  हे $\delta_1$
  चे उघडे गृहीतक असेल,
  तर \Intro{\lif}
  अनुमानात ते मुक्त होते,
  आणि~$\delta$ ची
  उघडी गृहीतके
  $\Gamma =
  \Gamma' \setminus \{x:B\}$
  होतात. म्हणजे
  ~$\delta_1'$ ही
  $x:B, \Gamma \Sequent C$
  ची सिद्धता
  असते; मग
  $\delta'$ पुढीलप्रमाणे
  घेता येते:
  \[
  \AxiomC{}
  \RightLabel{$\delta_1'$}
  \Deduce$x:B, \Gamma \fCenter C$
  \RightLabel{\Intro{\lif}}
  \UnaryInf$\Gamma \fCenter B \lif C$
  \DisplayProof
  \]
  $B^x$ हे~$\delta_1$
  चे उघडे गृहीतक
  नसेल, तर
  \Intro{\lif}
  कोणतेही गृहीतक
  मुक्त करत नाही;
  ~$\delta$ आणि
  $\delta_1$ यांची
  उघडी गृहीतके तीच
  राहतात. ~\Log{N2c}
  आणि~\Log{N2i}
  मध्ये \Intro{\lif}
  च्या आधारविधानाच्या
  संदर्भात $x:B$
  असणे आवश्यक नसल्याने
  $\delta'$ पुढीलप्रमाणे
  घेता येते:
  \[
  \AxiomC{}
  \RightLabel{$\delta_1'$}
  \Deduce$\Gamma \fCenter C$
  \RightLabel{\Intro{\lif}}
  \UnaryInf$\Gamma \fCenter B \lif C$
  \DisplayProof
  \]

  ~$\delta$ चा शेवट
  इतर नियमांत होणारे
  प्रसंगही असेच आहेत.
\end{proof}

\begin{prop}
$\delta$ ही \Log{N2c}
किंवा \Log{N2i} मध्ये
$\Gamma \Sequent A$
ची सिद्धता असेल,
तर~\Log{N1c}
(\Log{N1i}) मध्ये
अंतिम-सूत्र~$A$
आणि प्रत्येक
$x:B \in \Gamma$
साठी उघडे गृहीतक~$B^x$
असलेली सिद्धता~$\delta'$
मिळते.
\end{prop}

\begin{proof}
  पुन्हा~$\delta$
  ची उंची~$n$
  यावर विगमन करू.
  $n=0$ असेल, तर
  $\delta$ हे
  $x:A \Sequent A$
  स्वयंसिद्धक असते.
  सिद्धता~$\delta'$
  म्हणजे गृहीतक~$A^x$.

  $n>0$ असेल,
  तर~$\delta$ मधील
  अंतिम अनुमानानुसार
  प्रसंग पाडू.
  उदाहरणार्थ, ते
  ~$\Intro{\lif}$
  असेल, तर~$\delta$
  चा शेवट असा:
  \[
  \bottomAlignProof
  \AxiomC{}
  \RightLabel{$\delta_1$}
  \Deduce$x:B, \Gamma \fCenter C$
  \RightLabel{\Intro{\lif}}
  \UnaryInf$\Gamma \fCenter B \lif C$
  \DisplayProof
  \quad\text{ किंवा }\quad
  \bottomAlignProof
  \AxiomC{}
  \RightLabel{$\delta_1$}
  \Deduce$\Gamma \fCenter C$
  \RightLabel{\Intro{\lif}}
  \UnaryInf$\Gamma \fCenter B \lif C$
  \DisplayProof
  \]
  विगमन गृहीतकानुसार
  ~\Log{N1c} (\Log{N1i})
  मध्ये~$C$ ची
  सिद्धता मिळते.
  पहिल्या प्रसंगी
  $B^x$ हे~$\delta_1'$
  चे उघडे गृहीतक असते;
  दुसऱ्या प्रसंगी नसते.
  मग सिद्धता~$\delta'$
  अनुक्रमे पुढीलप्रमाणे
  घेता येते:
  \[
  \bottomAlignProof
  \AxiomC{$\Discharge{B}{x}$}
  \RightLabel{$\delta_1'$}
  \DeduceC{$C$}
  \DischargeRule{\Intro{\lif}}{x}
  \UnaryInfC{$B \lif C$}
  \DisplayProof
  \quad\text{ किंवा }\quad
  \bottomAlignProof
  \AxiomC{}
  \RightLabel{$\delta_1'$}
  \DeduceC{$C$}
  \RightLabel{\Intro{\lif}}
  \UnaryInfC{$B \lif C$}
  \DisplayProof
  \]
\end{proof}













\section{नियमित सिद्धता आणि आदेशन}\label{pt:nat:qua:sec}

संख्यापकांच्या नियमांमध्ये
~\Elim{\lexists} आणि
~\Intro{\lforall} यांना
लागणाऱ्या “आयगेन चराच्या
अटींमुळे” काही गुंतागुंत येते.
या अनुमानांतील
स्थिरांक~$c$ पुन्हा
वापरला जाणार नाही याची
खात्री केल्यास बहुतांश
गुंतागुंत हाताळता येते.
पुढील व्याख्या करू.

\begin{defn}
नैसर्गिक निगमनातील
सिद्धता~$\delta$
\emph{नियमित} असते, जर
\begin{enumerate}
  \item कोणतेही आयगेन
  चर~$a$ एकाहून अधिक
  \Elim{\lexists} किंवा
  ~\Intro{\lforall} अनुमानांचे
  आयगेन चर म्हणून वापरलेले
  नसेल; आणि
  \item $a$ हे
  \Elim{\lexists} किंवा
  ~\Intro{\lforall} अनुमानाचे
  आयगेन चर असेल, तर ते
  अनुक्रमे गौण आधारविधानाकडे
  किंवा आधारविधानाकडे
  जाणाऱ्या निकटच्या
  उप-सिद्धता मध्येच
  येत असेल.
\end{enumerate}
\end{defn}

\begin{lem}\label{pt:nat:qua:lem:subst-regular}
$\delta$ ही $C$ मध्ये
संपणारी नियमित सिद्धता असू द्या.
$c$ हा~$\delta$ मध्ये
आयगेन चर म्हणून न वापरलेला
स्थिरांक असू द्या;
आणि~$t$ हे $\delta$
चे कोणतेही आयगेन चर नसलेले
बंद पद असू द्या. या प्रकरणातील संख्यापक-नियमांत
वापरलेल्या बंद-पदाच्या अटीनुसार ही मर्यादा ठेवली आहे. मग~$\delta$
मधील $c$ च्या प्रत्येक
उपस्थितीच्या जागी $t$
बसवून मिळणारी
$\Subst{\delta}{t}{c}$
ही नियमित सिद्धता आहे.
$\Subst{\delta}{t}{c}$
चे अंतिम-सूत्र
$\Subst{C}{t}{c}$ असते.
$A^x$ हे~$\delta$
चे उघडे गृहीतक असेल,
तर $\Subst{A}{t}{c}^x$
हे $\Subst{\delta}{t}{c}$
चे उघडे गृहीतक असते.
\end{lem}

\begin{proof}
  $\pheight{\delta}$ वर
  विगमन करू.
  $\pheight{\delta} = 0$
  असेल, तर सिद्धतेत केवळ
  गृहीतक~$A^x$ असते.
  म्हणून $\Subst{\delta}{t}{c}$
  हीच~$\Subst{A}{t}{c}^x$
  असते.

  $\pheight{\delta} > 0$
  असेल, तर~$\delta$
  मधील अंतिम अनुमानानुसार
  प्रसंग पाडू.
  विधानीय संयोजकांच्या
  नियमांचे प्रसंग सोपे
  असून सरावासाठी ठेवले आहेत.
  $\delta$ चा शेवट
  संख्यापक अनुमानात होतो
  ते प्रसंग पाहू.
  \begin{enumerate}
    \item अंतिम अनुमान
    $\Intro{\lforall}$
    असेल, तर~$\delta$
    चे रूप असे आहे:
    \[
    \AxiomC{}
    \RightLabel{$\delta'$}
    \DeduceC{$A(a,c)$}
    \RightLabel{\Intro{\lforall}}
    \UnaryInfC{$\lforall[x][A(x,c)]$}
    \DisplayProof
    \]
    विगमन गृहीतकानुसार
    $\Subst{\delta'}{t}{c}$
    ही $A(a,t)$ ची
    नियमित सिद्धता आहे.
    $a$ हे~$\delta$
    चे आयगेन चर असल्याने
    $a$ हे~$t$ मध्ये येत नाही.
    म्हणून
    $\Subst{\delta}{t}{c}$
    मधील अंतिम अनुमान
    \[
    \AxiomC{}
    \RightLabel{$\Subst{\delta'}{t}{c}$}
    \DeduceC{$A(a,t)$}
    \RightLabel{\Intro{\lforall}}
    \UnaryInfC{$\lforall[x][A(x,t)]$}
    \DisplayProof
    \]
    योग्य आहे: $t$
    मध्ये~$a$ नसल्याने
    निष्कर्षात किंवा कोणत्याही
    उघड्या गृहीतकात~$a$
    येत नाही.
    \item अंतिम अनुमान
    $\Elim{\lforall}$
    असेल, तर~$\delta$
    चे रूप असे आहे:
    \[
    \AxiomC{}
    \RightLabel{$\delta'$}
    \DeduceC{$\lforall[x][A(x,c)]$}
    \RightLabel{\Elim{\lforall}}
    \UnaryInfC{$A(s(c),c)$}
    \DisplayProof
    \]
    विगमन गृहीतकानुसार
    $\Subst{\delta'}{t}{c}$
    ही $\lforall[x][A(x,t)]$
    ची नियमित सिद्धता
    आहे. (निष्कर्षातील
    पद~$s(c)$ मध्ये~$c$
    असूही शकतो किंवा नसूही
    शकतो.) आता
    $\Subst{\delta}{t}{c}$
    मधील अंतिम अनुमान
    \[
    \AxiomC{}
    \RightLabel{$\Subst{\delta'}{t}{c}$}
    \DeduceC{$\lforall[x][A(x,t)]$}
    \RightLabel{\Elim{\lforall}}
    \UnaryInfC{$A(s(t),t)$}
    \DisplayProof
    \]
    योग्य आहे, आणि
    $A(s(t),t) \ident
    \Subst{A(s(c),c)}{t}{c}$.
    \Intro{\lexists} चा
    प्रसंगही असाच आहे.
  \item अंतिम अनुमान
    $\Elim{\lexists}$
    असेल, तर~$\delta$
    चे रूप असे आहे:
    \[
    \AxiomC{}
    \RightLabel{$\delta_1'$}
    \DeduceC{$\lexists[x][A(x,c)]$}
    \AxiomC{$\Discharge{A(a,c)}{x}$}
    \RightLabel{$\delta_2'$}
    \DeduceC{$C$}
    \DischargeRule{\Elim{\lexists}}{x}
    \BinaryInfC{$C$}
    \DisplayProof
    \]
    विगमन गृहीतकानुसार
    $\Subst{\delta_1'}{t}{c}$
    आणि
    $\Subst{\delta_2'}{t}{c}$
    या अनुक्रमे
    $\lexists[x][A(x,t)]$
    ची आणि उघड्या
    गृहीतक~$A(a,t)$ पासून
    $\Subst{C}{t}{c}$
    ची नियमित सिद्धता
    आहेत. $a$ हे~$\delta$
    चे आयगेन चर असल्याने
    $a$ हे~$t$ मध्ये येत नाही.
    म्हणून
    $\Subst{\delta}{t}{c}$
    मधील अंतिम अनुमान
    \[
    \AxiomC{}
    \RightLabel{$\Subst{\delta_1'}{t}{c}$}
    \DeduceC{$\lexists[x][A(x,t)]$}
    \AxiomC{$\Discharge{A(a,t)}{x}$}
    \RightLabel{$\Subst{\delta_2'}{t}{c}$}
    \DeduceC{$\Subst{C}{t}{c}$}
    \DischargeRule{\Elim{\lexists}}{x}
    \BinaryInfC{$\Subst{C}{t}{c}$}
    \DisplayProof
    \]
    योग्य आहे:
    $\Subst{A(x,t)}{a}{x}
    \ident \Subst{A(a,c)}{t}{c}$;
    निष्कर्ष~$\Subst{C}{t}{c}$
    किंवा कोणत्याही उघड्या
    गृहीतकात~$a$ येत नाही.
  \end{enumerate}
\end{proof}

\begin{prop}\label{pt:nat:qua:prop:regular}
$\delta$ ही \Log{N1c}
किंवा \Log{N1i} मधील
सिद्धता असेल,
तर त्याच रचनेची
(विशेषतः त्याच
उंची ची)
नियमित सिद्धता~$\delta'$
असते; ती~$\delta$
पेक्षा आयगेन चरांची आणि आवश्यक असल्यास
बद्ध गृहीतक-खुणांची नावे बदलल्यामुळे भिन्न असते;
उघडी गृहीतके आणि त्यांची नावे जपली जातात.
\end{prop}

\begin{proof}
या सिद्धतेपुरते,
~$\delta$ मधील एखादे
आयगेन-चर अनुमान
\emph{स्वच्छ} म्हणू,
जर त्याचे आयगेन चर
इतर कोणत्याही अनुमानाचे
आयगेन चर नसेल आणि
त्या अनुमानाच्या निकटच्या
उप-सिद्धता व्यतिरिक्त
~$\delta$ मध्ये
कुठेही येत नसेल;
अन्यथा ते \emph{अस्वच्छ}
म्हणू. ~$n$ ही~$\delta$
मधील अस्वच्छ आयगेन-चर
अनुमानांची संख्या असू द्या.
~$n$ वर विगमन करून
~$\delta$ चे अपेक्षित
नियमित सिद्धता~$\delta'$
मध्ये रूपांतर करता येते,
हे दाखवू. $n=0$ असेल,
तर~$\delta$ मधील सर्व
आयगेन-चर अनुमाने स्वच्छ
आहेत; म्हणून~$\delta$
नियमित आहे आणि आणखी
काही सिद्ध करायचे नाही.
अन्यथा सर्वात वरचे
अस्वच्छ आयगेन-चर अनुमान
निवडा; उदाहरणार्थ
\Intro{\lforall}.
त्यात संपणारी~$\delta$
ची उप-सिद्धता~$\delta_1$
या रूपाची आहे:
    \[
    \AxiomC{}
    \RightLabel{$\delta_2$}
    \DeduceC{$A(c)$}
    \RightLabel{\Intro{\lforall}}
    \UnaryInfC{$\lforall[x][A(x)]$}
    \DisplayProof
    \]
~$c$ हे आयगेन चर असल्याने
ते~$\lforall[x][A(x)]$
मध्ये किंवा उघड्या गृहीतकात
येत नाही. ~$\delta_2$
नियमित आहे: निवडलेल्या
अनुमानाच्या वर असलेली
सर्व आयगेन-चर अनुमाने
स्वच्छ आहेत. ~$a$ हा
~$\delta$ मध्ये कुठेही
न येणारा स्थिरांक
असू द्या. आदेशनापूर्वी $\delta_2$ मधील
बद्ध गृहीतक-खुणा व त्यांना बांधलेली गृहीतके एकत्र
नवी नावे देऊन उर्वरित $\delta$ च्या खुणांपासून
वेगळी करा; उघड्या गृहीतकांच्या खुणा बदलू नका.
त्यामुळे आदेशनाने बदललेले एखादे अंतर्गत गृहीतक
बाहेरच्या वेगळ्या सूत्राबरोबर तीच खूण वापरणार नाही.
\ref{pt:nat:qua:lem:subst-regular}
नुसार
$\Subst{\delta_2}{a}{c}$
ही $A(a)$ ची नियमित
सिद्धता आहे.
आयगेन चराच्या अटीने
~$\delta_2$ च्या कोणत्याही
उघड्या गृहीतकात~$c$
नसतो; म्हणून
~$\Subst{\delta_2}{a}{c}$
च्या उघड्या गृहीतकात
~$a$ नसतो. त्यामुळे
~$\Intro{\lforall}$
लावता येतो.
~$\delta$ मधील
$\delta_1$ च्या जागी
पुढील सिद्धता~$\delta_1'$
ठेवा:
    \[
    \AxiomC{}
    \RightLabel{$\Subst{\delta_2}{a}{c}$}
    \DeduceC{$A(a)$}
    \RightLabel{\Intro{\lforall}}
    \UnaryInfC{$\lforall[x][A(x)]$}
    \DisplayProof
    \]
यात एक अस्वच्छ आयगेन-चर
अनुमान स्वच्छ केले आहे,
आयगेन चर~$c$ च्या जागी~$a$ आले आहे,
आणि आवश्यक बद्ध गृहीतक-खुणा नवी नावे देऊन जपल्या आहेत.
परिणामी सिद्धतेत
जास्तीत जास्त $n-1$ अस्वच्छ आयगेन-चर
अनुमाने राहतात; जुने चर इतरत्र न उरल्यास
आणखी अनुमानेही स्वच्छ होऊ शकतात. म्हणून
विगमन गृहीतकाने अपेक्षित
निष्कर्ष मिळतो.
\end{proof}

\begin{prob}
\ref{pt:nat:qua:prop:regular} च्या
सिद्धतेत सर्वात वरचे
आयगेन-चर अनुमान
~\Elim{\lexists} असते
तो प्रसंग वर्णन करा.
\end{prob}

नुकतेच सिद्ध केलेले
विधान आपण पुढे स्पष्टपणे
वापरणार नाही; पण चर्चेतील
सर्व सिद्धता नियमित
आहेत असे मानू.
तसे नसल्यास आयगेन
चरांची नावे बदलून त्या
नियमित करता येतात.

\ref{pt:nat:qua:lem:subst-regular}
आणि \ref{pt:nat:qua:prop:regular}
ही विधाने \Log{N2c}
आणि~\Log{N2i}
यांनाही लागू होतात.
त्यांच्या सिद्धता
सरावासाठी ठेवतो.













\section{सिद्धता यांचे कलम-संयोजन}\label{pt:nat:gra:sec}

~\Log{N1c} आणि~\Log{N1i}
मधील सिद्धता म्हणजे
गृहीतके~$B$ यांपासून
सूत्रे~$A$ यांच्या
सिद्धता होत. आता $B$
चीही सिद्धता आहे असे समजा.
$B$ ची सिद्धता आणि
$B$ पासून $A$ ची सिद्धता
एकत्र करून, $B$ हे गृहीतक
न वापरणारी $A$ ची
सिद्धता कशी मिळवता येईल?

एक मार्ग म्हणजे
~\Intro{\lif} लावून $A$
च्या सिद्धता मधील
$B$ हे गृहीतक काढून टाकणे.
मग मिळालेली $B \lif A$
ची सिद्धता आणि $B$
ची सिद्धता एकत्र
केल्यास पुढील सिद्धता मिळते:
\[
\AxiomC{$\Discharge{B}{x}$}
\DeduceC{$A$}
\DischargeRule{\Intro{\lif}}{x}
\UnaryInfC{$B \lif A$}
\AxiomC{}
\DeduceC{$B$}
\RightLabel{\Elim{\lif}}
\BinaryInfC{$A$}
\DisplayProof
\]
हे सरळ आहे, पण काहीसे
अनावश्यक वळण घेते:
$\lif$ आणून लगेच त्याचे
निराकरण का करावे?
$B \lif A$ मधून जाणारा
असा “वळसा” टाळता यायला हवा.
(असे वळसे नेहमी टाळता
येतात, हाच प्रत्यक्षात
सामान्यरूपीकरण प्रमेयाचा आशय आहे.)

~\Log{N1c} आणि~\Log{N1i}
मध्ये या प्रसंगी तो वळसा
सहज टाळता येतो: $B^x$
या गृहीतकांच्या जागी $B$
ची संपूर्ण सिद्धता बसवता येते.
इतर सिद्धतांच्या गृहीतकांवर
सिद्धता चे असे
“कलम-संयोजन” उपयुक्त आहे;
म्हणून त्याची व्याख्या नोंदवू.

\begin{defn}
समजा $\delta_1$ ही उघडे
गृहीतक~$B^x$ यापासून
$A$ ची \Log{N1c}-सिद्धता
(किंवा \Log{N1i}-सिद्धता) आहे,
आणि $\delta_2$ ही $B$
ची \Log{N1c}-सिद्धता
(\Log{N1i}-सिद्धता) आहे.
मग $\delta_1$ मधील
$B^x$ च्या प्रत्येक
निरसन न झालेल्या उदाहरणाच्या
जागी $\delta_2$ बसवून
मिळणारा परिणाम म्हणजे
$\Subst{\delta_1}{\delta_2}{B^x}$.
\end{defn}

~$\delta_1$ आणि $\delta_2$
मध्ये एकाच खूणेची, पण वेगवेगळी
गृहीतक-सूत्रे नसतील,
आणि $\delta_2$ च्या कोणत्याही
उघड्या गृहीतकात $\delta_1$
चे आयगेन चर येत नसेल,
आणि $\delta_2$ च्या उघड्या गृहीतकाची कोणतीही खूण
$\delta_1$ मधील मुक्ती-नियमाची खूण नसेल,
तर परिणामी
$\Subst{\delta_1}{\delta_2}{B^x}$
ही योग्य सिद्धता आहे.
तिची उघडी गृहीतके म्हणजे
$\delta_1$ मधील बसवलेल्या
उदाहरणांव्यतिरिक्त
उरलेली उघडी गृहीतके,
आणि $\delta_2$ मधील
उघडी गृहीतके.
सिद्धता चे कलम-संयोजन
करताना या अटी सहसा पूर्ण होतात.
~$\delta_1$ आणि $\delta_2$
या मोठ्या सिद्धता
च्या उपसिद्धता असतील,
तर पहिली अट पूर्ण होते.
दुसरी अट पूर्ण नसेल,
तर $\delta_1$ ची आयगेन चरे
पुन्हा नाव देऊन ती पूर्ण
करता येते.
ही खुणांची नवी अटही आवश्यक आहे: दोन्हीकडे तेच सूत्र
असले तरी आत बसवलेले उघडे गृहीतक बाहेरच्या मुक्ती-नियमाने
अनवधानाने मुक्त होता कामा नये. कलम-संयोजनापूर्वी
$\delta_1$ मधील मुक्ती-खुणा आणि त्यांना बांधलेली गृहीतके
एकत्र नवी नावे देऊन, $\delta_2$ च्या खुणांपासून वेगळी
करता येतात; मूळ उघड्या गृहीतकांची नावे जपायची.

तशीच व्याख्या
\Log{N2c} आणि~\Log{N2i}
यांनाही लागू होते.

\begin{defn}
समजा $\delta_1$ ही
$x:B, \Gamma_1 \Sequent A$
ची \Log{N2c}-सिद्धता
(किंवा \Log{N2i}-सिद्धता) आहे;
आणि $\delta_2$ ही
$\Gamma_2 \Sequent B$
ची \Log{N2c}-सिद्धता
(\Log{N2i}-सिद्धता) आहे.
मग $\delta_1$ मधील
अंतिम क्रमवर्तीपर्यंत मुक्त न झालेल्या
$x:B$ शी संबंधित प्रत्येक $x:B \Sequent B$
स्वयंसिद्धकाच्या जागी $\delta_2$ बसवून,
त्याखालील प्रत्येक क्रमवर्तीच्या संदर्भातील
संबंधित $x:B$ काढून त्याऐवजी $\Gamma_2$ जोडून
मिळणारा परिणाम म्हणजे
$\Subst{\delta_1}{\delta_2}{x:B}$.
\end{defn}

\begin{prop}
जर $\delta_1 \Proves
x:B, \Gamma_1 \Sequent
A$ आणि $\delta_2 \Proves
\Gamma_2 \Sequent B$,
तर $\Subst{\delta_1}{\delta_2}{x:B}
\Proves \Gamma_1, \Gamma_2 \Sequent
A$, मात्र दोन्ही सिद्धतांतील गृहीतक-खुणा सुसंगत असाव्यात,
$\Gamma_2$ मधील कोणतीही खूण $\delta_1$ मधील
मुक्ती-नियमाची खूण नसावी आणि $\delta_1$ चे
कोणतेही आयगेन चर $\Gamma_2$ मध्ये येता कामा नये.
\end{prop}

\begin{proof}
  $\pheight{\delta_1}$ वर
  विगमन करून.
\end{proof}












\section{\Log{G2i} पासून \Log{N2i} मध्ये रूपांतर}\label{pt:nat:tgi:sec}

~\Log{G2i} मधील क्रमवर्तींचे पूर्वपक्ष हे
सूत्रे यांचे बहुसंच~$\Gamma$ असतात; तर
\Log{N2i} मधील क्रमवर्तींचे पूर्वपक्ष हे
खूण लावलेल्या सूत्रांचे संच~$\Gamma'$ असतात.
त्यांतील प्रत्येक खूण फक्त एकदाच येते.
खूण लावलेल्या सूत्रांचा संच~$\Gamma'$
हा बहुसंच~$\Gamma$ शी \emph{अनुरूप आहे} असे
तेव्हाच म्हणू, जेव्हा प्रत्येक सूत्र~$A$
साठी $x : A$ या रूपातील~$\Gamma'$ चे
घटक जितके आहेत तितकी संख्या
~$\Gamma$ मध्ये~$A$ जितक्या वेळा येते
(म्हणजे~$A$ च्या प्रती जितक्या आहेत) त्या
संख्येपेक्षा अधिक नसते.

\begin{prop}\label{pt:nat:tgi:prop:G2i-to-N2i}
$\Log{G2i} + \Cut \Proves \Gamma \Sequent \Delta$
असेल, तर
$\Log{N2i} \Proves \Gamma' \Sequent C$.
येथे~$\Delta$ रिकामा असेल तर
$C \ident \lfalse$; अन्यथा
$\Delta = \{A\}$ आणि $C \ident A$.
तसेच~$\Gamma'$ हा~$\Gamma$ शी अनुरूप आहे.
\end{prop}

\begin{proof}
  ~$\pi$ ही~\Log{G2i} + \Cut मध्ये
  $\Gamma \Sequent \Delta$ ची
  सिद्धता असेल, तर~$\Gamma$
  शी अनुरूप असा~$\Gamma'$ आणि
  ~\Log{N2i} मध्ये
  ~$\Gamma' \Sequent C$ ची
  सिद्धता~$\delta$ मिळते,
  हे दाखवू. त्यासाठी
  $n = \pheight{\pi}$
  यावर विगमन करू.

  $n = 0$ असेल, तर~$\pi$
  हे स्वयंसिद्धक आहे.
  ते $\lfalse \Sequent \quad$
  असे असेल, तर~$\delta$
  हे $x:\lfalse \Sequent \lfalse$
  स्वयंसिद्धक; म्हणजे
  $\Gamma' = \{x:\lfalse\}$
  आणि $C \ident \lfalse$.
  ~$\pi$ हे $A \Sequent A$
  या रूपाचे स्वयंसिद्धक असेल,
  तर~$\delta$ हे
  $x: A \Sequent A$.

  $n>0$ असेल, तर~$\pi$
  चा शेवट~$R$ या नियमात होतो.
  विगमन गृहीतकानुसार त्या नियमाच्या
  आधारविधानांशी अनुरूप क्रमवर्तींच्या
  सिद्धता ~\Log{N2i} मध्ये मिळतात.
  या सिद्धतांतील सर्व मुक्तीच्या खुणा
  वेगवेगळ्या आहेत, आणि एखादी खूण~$x$
  मुक्त होत असेल तर तिला मुक्त करणाऱ्या
  अनुमानाच्या वरच येते, असे गृहीत धरू
  शकतो; गरज पडल्यास खुणा बदलू शकतो.
  या सिद्धता एकत्र करून योग्य
  $\Gamma' \Sequent C$ ची सिद्धता
  कशी मिळते ते आता~$R$ नुसार पाहू.

  \begin{enumerate}
    \item $R$ हा \RightR{\Weakening} असेल,
    तर~$\pi$ चा शेवट असा होतो:
    \[
    \AxiomC{}
    \RightLabel{$\pi_1$}
    \Deduce$\Gamma \fCenter$
    \RightLabel{\RightR{\Weakening}}
    \UnaryInf$\Gamma \fCenter A$
    \DisplayProof
    \]
    ~$\pi_1$ वर विगमन गृहीतक लावल्यास
    $\Gamma' \Sequent \lfalse$ ची
    सिद्धता मिळते. तिला~$\FalseInt$
    लावून $\Gamma' \Sequent A$ ची
    सिद्धता मिळते.
    \item $R$ हा \LeftR{\Weakening} असेल,
    तर~$\pi$ चा शेवट असा होतो:
    \[
    \AxiomC{}
    \RightLabel{$\pi_1$}
    \Deduce$\Gamma \fCenter \Delta$
    \RightLabel{\LeftR{\Weakening}}
    \UnaryInf$A, \Gamma \fCenter \Delta$
    \DisplayProof
    \]
    ~$\pi_1$ वर विगमन गृहीतक लावल्यास
    $\Gamma' \Sequent C$ ची
    सिद्धता मिळते; तीच~$\delta$ घेऊ.
    ~$\Gamma'$ हा~$\Gamma$ शी अनुरूप असेल,
    तर तो $A, \Gamma$ शीही अनुरूप असतो.
    कारण~$\Gamma'$ मधील $x : A$ ची संख्या
    ही~$\Gamma$ मधील~$A$ च्या प्रतींपेक्षा
    अधिक नाही; त्यामुळे ती
    $A, \Gamma$ मधील~$A$
    च्या प्रतींपेक्षाही अधिक नाही.
    \item $R$ हा \LeftR{\Contraction} असेल,
    तर~$\pi$ चा शेवट असा होतो:
    \[
    \AxiomC{}
    \RightLabel{$\pi_1$}
    \Deduce$A, A, \Gamma \fCenter \Delta$
    \RightLabel{\LeftR{\Contraction}}
    \UnaryInf$A, \Gamma \fCenter \Delta$
    \DisplayProof
    \]
    ~$\pi_1$ वर विगमन गृहीतक लावल्यास
    $\Pi, \Gamma' \Sequent C$ ची
    सिद्धता~$\delta_1$ मिळते;
    येथे $\Pi \subseteq \{x : A,
    y:A\}$. $\Pi = \{x : A, y:A\}$
    असेल, तर~$\delta_1$ मधील
    सर्व खूण लावलेली सूत्रे~$y:A$
    ही~$x:A$ ने बदला.
    ~$\delta_1$ च्या अंतिम क्रमवर्तीच्या
    संदर्भात~$x$ असल्यामुळे
    ~$\delta_1$ मधील कोणतेही अनुमान
    $x:B$ या रूपाचे सूत्र
    मुक्त करत नाही, असे धरता येते;
    म्हणजे~$\delta_1$ मधील
    कोणत्याही क्रमवर्तीच्या डाव्या बाजूवरील
    $x:B$ सक्रिय नसते.
    म्हणून बदललेली सिद्धता अजूनही
    ~\Log{N2i} मधील सिद्धता असते.
    \item $R$ हा $\RightR{\lif}$ असेल,
    तर~$\pi$ चा शेवट असा होतो:
    \[
    \AxiomC{}
    \RightLabel{$\pi_1$}
    \Deduce$A, \Gamma \fCenter B$
    \RightLabel{\RightR{\lif}}
    \UnaryInf$\Gamma \fCenter A \lif B$
    \DisplayProof
    \]
    विगमन गृहीतकानुसार
    $x: A, \Gamma' \Sequent B$
    किंवा $\Gamma' \Sequent B$
    ची सिद्धता~$\delta_1$ मिळते;
    येथे~$\Gamma'$ हा~$\Gamma$ शी
    अनुरूप आहे. ~$\delta_1$ ला
    \Intro{\lif} चे एक अनुमान जोडून
    मिळणारी सिद्धता~$\delta$ घेऊ.
    \item $R$ हा \LeftR{\lif} असेल,
    तर~$\pi$ चा शेवट असा होतो:
    \[
    \AxiomC{}
    \RightLabel{$\pi_1$}
    \Deduce$\Gamma_1 \fCenter A$
    \AxiomC{}
    \RightLabel{$\pi_2$}
    \Deduce$B, \Gamma_2 \fCenter \Delta$
    \RightLabel{\LeftR{\lif}}
    \BinaryInf$A \lif B, \Gamma_1, \Gamma_2 \fCenter \Delta$
    \DisplayProof
    \]
    विगमन गृहीतकानुसार सिद्धता मिळतात:
    $\delta_1$ ही $\Gamma_1' \Sequent A$ ची;
    आणि~$\delta_2$ ही
    $x: B, \Gamma_2' \Sequent C$
    किंवा $\Gamma_2' \Sequent C$ ची.
    ~$\delta_2$ दुसऱ्या रूपात असेल,
    तर~$\delta_2$ हीच~$\delta$ घेऊ.
    अन्यथा~$\delta_1$ मधील
    सर्व सूत्रांच्या आणि मुक्तीच्या खुणा
    बदलून~$\delta_1$ व~$\delta_2$
    यांत कोणतीही खूण समान नसल्याची
    खात्री करू शकतो. मग
    $\Gamma_1' \cup \Gamma_2'$ हा
    $\Gamma_1 \cup \Gamma_2$ शी अनुरूप आहे.
    ~$\delta_2$ च्या अंतिम क्रमवर्तीत
    $x:B$ असल्यामुळे
    ~$\delta_2$ मधील~$x$ या
    मुक्ती-खुणेच्या कोणत्याही अनुमानात
    $x:B$ सक्रिय नसते.
    ~$\delta_2$ मधील प्रत्येक
    $x:B \Sequent B$ स्वयंसिद्धकाच्या
    जागी पुढील सिद्धता~$\delta_3$ ठेवा:
    \[
    \Axiom$x: A \lif B \fCenter A \lif B$
    \AxiomC{}
    \RightLabel{$\delta_1$}
    \Deduce$\Gamma_1' \fCenter A$
    \RightLabel{\Elim{\lif}}
    \BinaryInf$x: A \lif B, \Gamma_1' \fCenter B$
    \DisplayProof
    \]
    आणि~$\delta_2$ मधील प्रत्येक
    $x:B$ च्या जागी $x:A \lif B$
    ठेवा. परिणामी
    $\Subst{\delta_2}{\delta_3}{x:B}$
    ही $x:A \lif B, \Gamma_1', \Gamma_2' \Sequent C$
    ची सिद्धता आहे.
    \item $R$ हा \RightR{\land} असेल,
    तर~$\pi$ चा शेवट असा होतो:
    \[
    \AxiomC{}
    \RightLabel{$\pi_1$}
    \Deduce$\Gamma_1 \fCenter A$
    \AxiomC{}
    \RightLabel{$\pi_2$}
    \Deduce$\Gamma_2 \fCenter B$
    \RightLabel{\RightR{\land}}
    \BinaryInf$\Gamma_1, \Gamma_2 \fCenter A \land B$
    \DisplayProof
    \]
    विगमन गृहीतकानुसार
    $\Gamma_1' \Sequent A$ ची
    सिद्धता~$\delta_1$ आणि
    $\Gamma_2' \Sequent B$ ची
    सिद्धता~$\delta_2$ मिळते.
    ~$\delta_2$ मधील सर्व सूत्रांच्या
    आणि मुक्तीच्या खुणा बदलून
    ~$\delta_1$ व~$\delta_2$ यांत
    कोणतीही खूण समान नसल्याची
    खात्री करता येते. मग
    $\Gamma_1' \cup \Gamma_2'$ हा
    $\Gamma_1 \cup \Gamma_2$ शी अनुरूप आहे.
    ~$\delta_1$ आणि~$\delta_2$
    यांच्या अंतिम क्रमवर्तींना
    $\Intro{\land}$ लावून
    सिद्धता~$\delta$ मिळते.
    \item $R$ हा $\LeftR{\land}$ असेल,
    तर, उदाहरणार्थ, ~$\pi$ चा
    शेवट असा होतो:
    \[
    \AxiomC{}
    \RightLabel{$\pi_1$}
    \Deduce$A, \Gamma \fCenter \Delta$
    \RightLabel{\LeftR{\land}}
    \UnaryInf$A \land B, \Gamma \fCenter \Delta$
    \DisplayProof
    \]
    विगमन गृहीतकानुसार
    $x: A, \Gamma' \Sequent C$
    किंवा $\Gamma' \Sequent C$
    ची सिद्धता~$\delta_1$ मिळते;
    येथे~$\Gamma'$ हा~$\Gamma$ शी
    अनुरूप आहे. दुसऱ्या प्रसंगी
    ~$\delta_1$ हीच~$\delta$ घेऊ.
    पहिल्या प्रसंगी~$x:A$
    अंतिम क्रमवर्तीत असल्यामुळे
    ~$\delta_1$ मधील~$x$ या
    मुक्ती-खुणेच्या कोणत्याही अनुमानात
    $x:A$ सक्रिय नसते.
    प्रत्येक $x:A \Sequent A$
    स्वयंसिद्धकाच्या जागी
    पुढील सिद्धता~$\delta_2$ ठेवा:
    \[
    \Axiom$x: A \land B \fCenter A \land B$
    \RightLabel{\Elim{\land}}
    \UnaryInf$x: A \land B \fCenter A$
    \DisplayProof
    \]
    आणि प्रत्येक~$x:A$ च्या जागी
    $x:A \land B$ ठेवा; म्हणजे
    $\Subst{\delta_1}{\delta_2}{x:A}$
    घ्या. ही $x:A \land B, \Gamma' \Sequent C$
    ची सिद्धता आहे.
    \item $R$ हा \Cut असेल,
    तर~$\pi$ चा शेवट असा होतो:
    \[
    \AxiomC{}
    \RightLabel{$\pi_1$}
    \Deduce$\Gamma_1 \fCenter A$
    \AxiomC{}
    \RightLabel{$\pi_2$}
    \Deduce$A, \Gamma_2 \fCenter \Delta$
    \RightLabel{\Cut}
    \BinaryInf$\Gamma_1, \Gamma_2 \fCenter \Delta$
    \DisplayProof
    \]
    विगमन गृहीतकानुसार
    $\Gamma_1' \Sequent A$ ची
    सिद्धता~$\delta_1$ आणि
    $x:A, \Gamma_2' \Sequent C$
    किंवा $\Gamma_2' \Sequent C$
    ची सिद्धता~$\delta_2$ मिळते.
    दुसऱ्या प्रसंगी~$\delta$
    म्हणून~$\delta_2$ घ्या.
    अन्यथा, गरज असल्यास
    ~$\delta_1$ मधील खुणा बदलून
    दोन सिद्धतांच्या खुणा वेगळ्या करा.
    ~$\delta_2$ मधील प्रत्येक
    $x:A \Sequent A$ स्वयंसिद्धकाच्या
    जागी~$\delta_1$ ठेवून
    $\Subst{\delta_2}{\delta_1}{x:A}$
    ही $\Gamma_1', \Gamma_2' \Sequent C$
    ची सिद्धता~$\delta$ मिळते.
  \end{enumerate}
\end{proof}













\section{\Log{N2} पासून \Log{G2} मध्ये रूपांतर}\label{pt:nat:tni:sec}

\begin{prop}\label{pt:nat:tni:prop:N2-to-G2}
$\Log{N2c}\ (\Log{N2i}) \Proves
\Gamma \Sequent A$ असेल, तर
$\Log{G2c}\ (\Log{G2i}) + \Cut \Proves \Gamma'
\Sequent \Delta$. येथे~$\Gamma$ मधील
खुणा काढून मिळणारा सूत्रांचा
बहुसंच~$\Gamma'$ आहे; तसेच~$A$
हा~$\lfalse$ नसेल तर
$\Delta = \{A\}$, आणि असेल तर
$\Delta = \emptyset$.
\end{prop}

\begin{proof}
पुन्हा~$\Gamma \Sequent A$ च्या
सिद्धता~$\delta$ च्या उंचीवर
विगमन करू. ~$\delta$ हे
$x:A \Sequent A$ स्वयंसिद्धक
असेल, तर~$\pi$ हे अनुरूप
$A \Sequent A$ स्वयंसिद्धक
आहे; किंवा~$A \ident \lfalse$
असेल तर
$\lfalse \Sequent \ $ स्वयंसिद्धक आहे.

~$\delta$ चा शेवट
अनुमाननियमात होत असेल, तर
पुढील प्रसंग वेगळे पाहू:
\begin{enumerate}
  \item $R$ हा
  मुख्य सूत्र~$B \lif C$
  असलेला~\Intro{\lif} नियम असेल,
  आणि~$x:B$ हे आधारविधानाच्या
  संदर्भात असेल पण निष्कर्षाच्या
  संदर्भात नसेल, तर~$\delta$
  चा शेवट असा होतो:
  \[
  \AxiomC{}
  \RightLabel{$\delta_1$}
  \Deduce$x:B, \Gamma \fCenter C$
  \RightLabel{\Intro{\lif}}
  \UnaryInf$\Gamma \fCenter B \lif C$
  \DisplayProof
  \]
  विगमन गृहीतकानुसार
  $B, \Gamma' \Sequent C$
  ची सिद्धता~$\pi_1$
  मिळते; $C \ident \lfalse$
  असेल तर तिचे रूप
  $B, \Gamma' \Sequent \ $
  असते. त्यावर~\RightR{\lif}
  लावून~$\pi$ मिळते;
  $C \ident \lfalse$ असताना
  त्याआधी~\RightR{\Weakening}
  लावावे लागते.
  \item $R$ हा~\Intro{\lif}
  असेल, पण~$x:B$
  आधारविधानाच्या संदर्भात नसेल,
  तर~$\delta$ चा शेवट असा होतो:
  \[
  \AxiomC{}
  \RightLabel{$\delta_1$}
  \Deduce$\Gamma \fCenter C$
  \RightLabel{\Intro{\lif}}
  \UnaryInf$\Gamma \fCenter B \lif C$
  \DisplayProof
  \]
  विगमन गृहीतकानुसार
  निष्पन्न सूत्र असत्य
  नसेल तेव्हा~$\Gamma' \Sequent C$
  ची सिद्धता~$\pi_1$ मिळते;
  असत्याच्या प्रसंगी तिचा
  उत्तरपक्ष रिकामा असतो.
  गरज असल्यास आधी
  उजवीकडील दुर्बलीकरणाने
  असत्य उजवीकडे आणा.
  मग~$B$ साठी
  \LeftR{\Weakening}
  आणि त्यानंतर
  \RightR{\lif} लावून~$\pi$
  मिळते.
  \item $R$ हा~\Elim{\lif}
  असेल, तर~$\delta$ चा
  शेवट असा होतो:
  \[
  \AxiomC{}
  \RightLabel{$\delta_1$}
  \Deduce$\Gamma_1 \fCenter B \lif A$
  \AxiomC{}
  \RightLabel{$\delta_2$}
  \Deduce$\Gamma_2 \fCenter B$
  \RightLabel{\Elim{\lif}}
  \BinaryInf$\Gamma_1, \Gamma_2 \fCenter A$
  \DisplayProof
  \]
  येथे~$\Gamma = \Gamma_1 \cup \Gamma_2$.
  विगमन गृहीतकानुसार
  सिद्धता मिळतात:
  $\pi_1$ ही
  $\Gamma_1' \Sequent B \lif A$
  ची आणि~$\pi_2$ ही
  $\Gamma_2' \Sequent B$ ची.
  दुसऱ्या निष्कर्षसूत्राचा असत्य
  हा प्रसंग असेल, तर दुसरी
  सिद्धता मिळवण्यासाठी
  विगमनातून आलेल्या रिकाम्या
  उत्तरपक्षाच्या सिद्धतेवर
  उजवीकडील दुर्बलीकरण लावा.
  पुढीलप्रमाणे सिद्धता~$\pi$
  मिळते:
  \[
  \AxiomC{}
  \RightLabel{$\pi_1$}
  \Deduce$\Gamma_1' \fCenter B \lif A$
  \AxiomC{}
  \RightLabel{$\pi_2$}
  \Deduce$\Gamma_2' \fCenter B$
  \Axiom$A \fCenter A$
  \RightLabel{\LeftR{\lif}}
  \BinaryInf$B \lif A, \Gamma_2' \fCenter A$
  \RightLabel{\Cut}
  \BinaryInf$\Gamma_1', \Gamma_2' \fCenter A$
  \DisplayProof
  \]
  बहुसंच~$\Gamma_1' \cup \Gamma_2'$
  हा $\Gamma' =
  (\Gamma_1 \cup \Gamma_2)'$
  याच्याशी समान असेलच असे नाही.
  उदाहरणार्थ,~$\Gamma_1$
  आणि~$\Gamma_2$ दोन्हींमध्ये
  $x:C$ असेल, तर
  $\Gamma_1' \cup \Gamma_2'$
  मध्ये~$C$ च्या दोन प्रती
  असतात; पण
  $(\Gamma_1 \cup \Gamma_2)'$
  मध्ये एकच प्रत असते.
  योग्य~\LeftR{\Contraction}
  अनुमाने जोडून हे दुरुस्त करता येते.

  $A \ident \lfalse$
  असेल, तर~$\delta_2$
  वरून विगमनाने मिळणारी
  सिद्धता एकटी रिकामा
  उत्तरपक्ष देत नाही.
  वरच्या रचनेत
  $A$ च्या ओळख-स्वयंसिद्धकाऐवजी
  रिकामा उत्तरपक्ष असलेले
  असत्य-स्वयंसिद्धक वापरा.
  डाव्या अभिव्यंजन-नियमापुढे
  कट लावल्यावर
  $\Gamma_1', \Gamma_2' \Sequent \ $
  ची सिद्धता~$\pi$ मिळते.
  त्यामुळे यासाठी
  \LeftR{\Weakening} किंवा
  \RightR{\Weakening}
  लावण्याची गरज नाही.
\item $R$ हा~\FalseCl असेल,
तर~$\delta$ चा शेवट असा:
  \[
  \AxiomC{}
  \RightLabel{$\delta_1$}
  \Deduce$x: \lnot A, \Gamma \fCenter \lfalse$
  \RightLabel{\FalseCl}
  \UnaryInf$\Gamma \fCenter A$
  \DisplayProof
  \]
  विगमन गृहीतकानुसार
  $\lnot A, \Gamma' \Sequent \ $
  ची सिद्धता~$\pi_1$
  आहे. पुढील सिद्धता~$\pi$ घेऊ:
  \[
  \Axiom$A \fCenter A$
  \RightLabel{\RightR{\lnot}}
  \UnaryInf$\fCenter A, \lnot A$
  \AxiomC{}
  \RightLabel{$\pi_1$}
  \Deduce$\lnot A, \Gamma' \fCenter $
  \RightLabel{\Cut}
  \BinaryInf$\Gamma' \fCenter A$
  \DisplayProof
  \]
  $A\ident\lfalse$ असेल, तर या शेवटच्या क्रमवर्तीचा
  उत्तरपक्षही रिकामा हवा. त्यासाठी वरील सिद्धतेला
  $\lfalse\Sequent\ $ या स्वयंसिद्धकासह आणखी एक
  $\Cut$ लावून $\Gamma'\Sequent\ $ मिळवायचे.
  $\FalseCl$ मध्ये कोणतेही गृहीतक मुक्त होत नसेल,
  तर आधाराची सिद्धता विगमनाने थेट
  $\Gamma'\Sequent\ $ देते; $A\not\ident\lfalse$
  असल्यासच उजवीकडील दुर्बलीकरणाने $A$ जोडायचे.
\end{enumerate}
फक्त~\FalseCl{} च्या
रूपांतरातून मिळणाऱ्या
सिद्धता~$\pi$ मध्ये
एखाद्या क्रमवर्तीच्या उजव्या बाजूस
एकापेक्षा अधिक सूत्र
येऊ शकते. म्हणजे
~\Log{N2i} मधील सिद्धता
रूपांतरित केल्यास
\Log{G2i} मधील सिद्धता मिळते.
\end{proof}



